\documentclass[11pt,a4paper]{article}

\usepackage[T1]{fontenc}
\usepackage[utf8]{inputenc}
\usepackage[english]{babel}
\frenchspacing
\usepackage{newtxtext,newtxmath}
\usepackage{xcolor}
\usepackage{graphicx}
\usepackage{booktabs}
\usepackage{enumitem}
\usepackage{microtype}
\usepackage[most]{tcolorbox}
\usepackage{titlesec}
\usepackage{fancyhdr}
\usepackage[a4paper,top=24mm,bottom=22mm,left=24mm,right=24mm,headsep=8mm]{geometry}
\usepackage{array}
\usepackage{siunitx}

% ---------------------------------------------------------------- UAI / SFU palette
\definecolor{uaired}{HTML}{A6192E}
\definecolor{sfuslate}{HTML}{54585A}
\definecolor{ink}{HTML}{1A1A1C}
\definecolor{muted}{HTML}{585C60}
\definecolor{faint}{HTML}{8B8F93}
\definecolor{rule}{HTML}{E4E2E2}
\definecolor{paper}{HTML}{F7F6F6}
\definecolor{good}{HTML}{1E6B3A}

\color{ink}

\newcommand{\gd}[1]{\textcolor{good}{#1}}
\newcommand{\bd}[1]{\textcolor{uaired}{#1}}

\newtcolorbox{keybox}{
  enhanced, breakable, colback=paper, colframe=rule,
  boxrule=0.5pt, arc=1.5pt, left=10pt, right=10pt, top=7pt, bottom=7pt, boxsep=0pt
}
\newtcolorbox{warnbox}{
  enhanced, breakable, colback=white, colframe=uaired,
  boxrule=0pt, leftrule=1.6pt, arc=0pt, outer arc=0pt,
  left=9pt, right=4pt, top=5pt, bottom=5pt, boxsep=0pt
}

\titleformat{\section}{\normalfont\large\bfseries\color{ink}}{\textcolor{uaired}{\thesection}}{0.7em}{}
\titleformat{\subsection}{\normalfont\normalsize\bfseries\color{sfuslate}}{\textcolor{sfuslate}{\thesubsection}}{0.6em}{}
\titlespacing*{\section}{0pt}{13pt}{6pt}
\titlespacing*{\subsection}{0pt}{10pt}{3pt}

\pagestyle{fancy}\fancyhf{}
\renewcommand{\headrulewidth}{0.4pt}
\renewcommand{\headrule}{\hbox to\headwidth{\color{rule}\leaders\hrule height \headrulewidth\hfill}}
\fancyhead[L]{\footnotesize\color{faint}\textsc{Parallel sweep --- performance campaign}}
\fancyhead[R]{\footnotesize\color{faint}UAI\,--\,SFU}
\fancyfoot[C]{\footnotesize\color{faint}\thepage}

\setlength{\parindent}{0pt}
\setlength{\parskip}{4pt}
\renewcommand{\arraystretch}{1.15}
\setlist[itemize]{leftmargin=14pt,itemsep=2pt,topsep=3pt}

\sisetup{group-separator={,},detect-weight=true}

\begin{document}

% ================================================================= TITLE BLOCK
\thispagestyle{empty}
\noindent
\begin{minipage}[c]{0.46\textwidth}\raggedright
  \includegraphics[height=12mm]{logo-uai.png}
\end{minipage}\hfill
\begin{minipage}[c]{0.46\textwidth}\raggedleft
  \includegraphics[height=12mm]{logo-sfu.png}
\end{minipage}

\vspace{14mm}

{\color{uaired}\rule{30mm}{2.2pt}}

\vspace{6mm}

{\fontsize{14}{17}\selectfont\color{muted}\textsc{Rolling-horizon feasible operating region}\par}

\vspace{3mm}

{\hyphenpenalty=10000\exhyphenpenalty=10000\raggedright
\fontsize{24}{29}\selectfont\bfseries
Parallel directional sweep\par}

\vspace{5mm}

{\fontsize{12.5}{17}\selectfont\color{muted}
Making the per-transformer rolling FOR \num{1.82}$\times$ faster:\\
what was changed, what it costs, and why the answers still hold\par}

\vspace{10mm}

{\footnotesize\color{faint}August 2026 \quad$\cdot$\quad Campaign report\par}

\vspace{8mm}

\begin{keybox}
\textbf{In one paragraph.} Each period of the rolling FOR solves 12 sweep
directions. They used to be solved one after another. They are now solved
several at a time, in separate CPLEX sessions. Nothing about the physics
changed: same exact quadratic form, same constraints, same switches. The
committed dispatch comes out \textbf{bit-identical}. Measured end to end on all
nine transformers under one controlled machine base, the campaign runs in
\textbf{19.6\,h} instead of \textbf{35.7\,h}.
\end{keybox}

\vspace{4mm}

% ================================================================= 1
\section{The headline}

\begin{center}
\begin{tabular}{@{}l r@{}}
\toprule
Whole per-transformer campaign  & \gd{\textbf{1.82$\times$}} \\
Best case (TR9, 90 batteries)   & 2.73$\times$ \\
Worst case (TR7, 259 batteries) & 1.31$\times$ \\
Solves                          & \gd{5184 of 5184 Optimal} \\
Committed dispatch (SOC files)  & \gd{bit-identical, all nine} \\
FOR boundary vertices           & within \SI{3}{\watt}--\SI{70}{\watt} \\
\bottomrule
\end{tabular}
\end{center}

The committed dispatch is the quantity the papers use. It is unchanged, exactly.
Only the FOR boundary vertices move, and by amounts far below the model's own
feasibility tolerance. Section~5 gives the numbers.

These are measurements, not projections. Every run was made one transformer per
fresh AIMMS session, on a machine held in a fixed state.

% ================================================================= 2
\section{How it works}

Within one period the 12 directions share the same feasible region. The data
reset is deterministic. The rolling map is fixed. Slot-1 dispatch is frozen.
And \texttt{FOR\_TieEps}~$=0$ makes each vertex independent of solve history.
So the directions may be solved in any order, or at the same time.

Each direction receives its own copy of the period's generated matrix and its
own solver session. Two settings control the split:

\begin{itemize}
  \item \texttt{RollAsyncK} --- how many directions solve at once (sessions);
  \item \texttt{RollAsyncThreads} --- threads per session.
\end{itemize}

The machine has 24 physical cores, so $K \times \text{threads} = 24$ is the budget.

Two derived quantities explain every result below:

\begin{itemize}
  \item \textbf{inflation} --- async mean solve divided by sequential mean solve:
        what concurrency costs each individual solve;
  \item \textbf{concurrency} --- total solver seconds divided by wall seconds:
        how many solves really overlapped.
\end{itemize}

Net speed-up is approximately concurrency divided by inflation. Neither number
alone predicts anything.

% ================================================================= 3
\section{Results --- all nine transformers}

\begin{center}
\begin{tabular}{@{}l r c r r r r r@{}}
\toprule
 & \textbf{batt} & \textbf{split} & \textbf{seq.} & \textbf{base} & \textbf{speed-up} & \textbf{infl.} & \textbf{conc.} \\
\midrule
TR4 &  46 & 6$\times$4 &  35.7 &  18.6 & 1.92$\times$ & 2.20 & 3.63 \\
TR3 &  49 & 6$\times$4 &  39.1 &  20.1 & 1.95$\times$ & 2.33 & 3.95 \\
TR1 &  77 & 6$\times$4 &  96.5 &  50.1 & 1.93$\times$ & 2.31 & 3.84 \\
TR9 &  90 & 6$\times$4 & 195.8 &  71.8 & \gd{\textbf{2.73$\times$}} & 1.59 & 3.78 \\
TR2 &  95 & 6$\times$4 & 140.4 &  71.4 & 1.97$\times$ & 2.21 & 3.79 \\
TR8 & 124 & 4$\times$6 & 298.9 & 125.1 & 2.39$\times$ & 1.18 & 2.42 \\
TR5 & 214 & 4$\times$6 & 516.5 & 204.4 & \gd{\textbf{2.53$\times$}} & 1.11 & 2.39 \\
TR7 & 259 & 4$\times$6 & 411.2 & 314.0 & \bd{1.31$\times$} & 1.97 & 2.17 \\
TR6 & 262 & 4$\times$6 & 409.1 & 301.6 & \bd{1.36$\times$} & 1.99 & 2.20 \\
\midrule
 & & & \textbf{2143.2} & \textbf{1177.0} & \textbf{1.82$\times$} & & \\
\bottomrule
\end{tabular}
\end{center}

Wall clock in minutes.

\subsection{The gain is not monotonic in size}

It peaks in the middle. TR5, at 214 batteries, beats every small transformer
except TR9; then the gain collapses above about 250. Three regimes, not a trend:

\begin{center}
\begin{tabular}{@{}p{0.30\textwidth} p{0.62\textwidth}@{}}
\toprule
\textbf{band} & \textbf{behaviour} \\
\midrule
below $\sim$120 batteries & $\sim$1.9$\times$; the split is \textbf{neutral} (Section~4) \\
124--214                  & \textbf{best} --- 2.4--2.5$\times$, inflation only 1.11--1.18 \\
above $\sim$250           & \textbf{ceiling $\sim$1.35$\times$} --- inflation $\sim$2.0 and no split fixes it \\
\bottomrule
\end{tabular}
\end{center}

TR7 and TR6 are the same result twice: 259 and 262 batteries, inflation 1.97 and
1.99, concurrency 2.17 and 2.20, speed-up 1.31$\times$ and 1.36$\times$. The
ceiling is structural, not a bad night. It is not memory exhaustion either ---
peak AIMMS \SI{14.8}{\giga\byte} with \SI{8.1}{\giga\byte} still free.

\subsection{Why TR9 tops the table}

TR9 and TR2 are nearly the same run: 71.8 against \SI{71.4}{\minute}, 28.29
against \SI{28.14}{\second} per solve, concurrency 3.78 against 3.79. The whole
difference in speed-up comes from the \emph{sequential} side, where TR9's solves
are intrinsically 40\,\% heavier (17.76 against \SI{12.73}{\second}). Heavier
solves carry the parallel overhead better.

\begin{warnbox}
Worth remembering when quoting these: the speed-up column is partly a statement
about the sequential baseline, not only about the parallel gain.
\end{warnbox}

% ================================================================= 4
\section{Choosing the split}

Three experiments. \textbf{TR4 tuning} (46 batteries, full runs): sequential
35m20s, 6$\times$2 28m16s, \textbf{6$\times$4 20m13s}, 12$\times$2 24m39s. Above
about six sessions the machine saturates.

\textbf{TR6 probe} (262 batteries, two periods per split, 24 threads throughout):

\begin{center}
\begin{tabular}{@{}l r r r@{}}
\toprule
\textbf{split} & \textbf{wall} & \textbf{mean solve} & \textbf{concurrency} \\
\midrule
6 $\times$ 4 & \SI{692}{\second} & \SI{95.4}{\second} & 3.31 \\
\textbf{4 $\times$ 6} & \gd{\textbf{621\,s}} & \SI{60.3}{\second} & 2.33 \\
3 $\times$ 8 & \SI{706}{\second} & \SI{53.4}{\second} & 1.82 \\
2 $\times$ 12 & \SI{746}{\second} & \SI{39.3}{\second} & 1.27 \\
\bottomrule
\end{tabular}
\end{center}

Fewer and fatter sessions cut the time per solve sharply, from \SI{95}{\second}
to \SI{39}{\second}, but the concurrency achieved is only $0.55$--$0.63 \times K$
whatever the split.

\textbf{TR3 split test} (49 batteries, full runs): 6$\times$4
\SI{20.05}{\minute}, 4$\times$6 \SI{20.02}{\minute}. A dead heat. The 4$\times$6
split cut the mean solve from 8.24 to \SI{5.44}{\second} exactly as predicted,
and concurrency fell from 3.95 to 2.61, cancelling it out.

\begin{keybox}
\textbf{Judge a split by wall clock, never by mean solve time.} On TR6,
2$\times$12 had the best mean solve of any split and the worst wall clock. On
TR3, 4$\times$6 improved the mean by a third and moved the wall by two seconds.
\end{keybox}

So the crossover is \emph{not} ``6$\times$4 is better below 120 batteries''. It is:

\begin{itemize}
  \item below $\sim$120 --- the split does not matter; 6$\times$4 is kept because
        there is no reason to change;
  \item above $\sim$120 --- 4$\times$6, and the margin is large;
  \item above $\sim$250 --- 4$\times$6 is still the best available, but nothing
        recovers the ceiling.
\end{itemize}

\subsection{How the probe held up}

The TR6 probe predicted inflation of 1.73 for 4$\times$6; the full 48-period run
gave 1.99, and the wall clock came in at \SI{5.03}{\hour} against a
\SI{4.5}{\hour} projection. Directionally right, quantitatively soft --- good
enough to pick a split, not to promise a wall clock.

% ================================================================= 5
\section{The batch wrapper cost 1.48$\times$ for nothing}

\texttt{RunFOR\_Rolling\_PerTrafo\_Async} used to load the whole feeder once at
the start, only to read the list of transformer tags. But
\texttt{RunFOR\_Rolling} empties everything and reloads as its first action. So
that load bought nothing, while it grew the AIMMS heap to full-feeder size for
the rest of the session.

\begin{center}
\begin{tabular}{@{}l r r r@{}}
\toprule
 & \textbf{standalone} & \textbf{in the batch} & \textbf{penalty} \\
\midrule
TR4 & \SI{7.76}{\second}/solve & \SI{11.58}{\second}/solve & 1.49$\times$ \\
TR6 & \SI{95.4}{\second}/solve & \SI{141.0}{\second}/solve & 1.48$\times$ \\
\bottomrule
\end{tabular}
\end{center}

The same penalty appears at different positions in the batch, so it was a
constant and not creeping fragmentation. Removing the load fixed it: TR4 went
from \SI{40.8}{\minute} to \SI{18.9}{\minute}.

This is why the campaign of 17--18 August is not quoted here. It ran across a
batch, across days and across machine loads, and produced two transformers
\emph{slower} than sequential. The table in Section~3 replaces it entirely.

% ================================================================= 6
\section{Is it still correct?}

Yes, with one caveat.

\begin{itemize}
  \item \textbf{Every solve Optimal} --- 5184 of 5184 across the nine runs.
  \item \textbf{Committed dispatch bit-identical} --- all nine SOC files match
        their sequential twins byte for byte.
  \item \textbf{Async is deterministic run to run.} Five transformers were
        compared against their August async runs and came back at exactly zero,
        including TR9 across two days and a reboot. The drift below is against
        the \emph{sequential} solver only.
  \item \textbf{FOR vertices drift a little, and the drift grows with size.}
\end{itemize}

\begin{center}
\begin{tabular}{@{}l r r@{}}
\toprule
 & \textbf{max $|\Delta \mathrm{proj}|$} & \textbf{rows over gate (of 576)} \\
\midrule
TR3, TR4      & \num{3e-6} p.u. (\SI{3}{\watt})    & 0 \\
TR1, TR2, TR9 & \num{7e-6}--\num{9e-6} p.u.        & 6--10 \\
TR8           & \num{2.2e-5} p.u.                  & 0 \\
TR5           & \num{3.9e-5} p.u.                  & 0 \\
TR6           & \num{4.8e-5} p.u.                  & 0 \\
TR7           & \num{7e-5} p.u. (\SI{70}{\watt})   & 1 \\
\bottomrule
\end{tabular}
\end{center}

The cause is the thread cap changing the order of arithmetic inside the barrier
solver. For scale, \SI{70}{\watt} is about \textbf{14 times inside} the model's
own \texttt{Feasibility\_Tolerance} of \num{1e-3} p.u.\ $=$ \SI{1}{\kilo\watt},
and some \num{0.005}\,\% of the extent of the region.

\begin{warnbox}
\textbf{Caveat.} If a number needs the strict \num{5e-6} gate, run that case
sequentially. The sequential \texttt{\_gmp} results stay the set to quote from.
\end{warnbox}

The invariant is \texttt{proj}, not $P$ and $Q$ on their own. Flat faces of the
region make the argmax a segment, so $P$ and $Q$ can slide while \texttt{proj}
holds.

% ================================================================= 7
\section{What to run}

\begin{center}
\begin{tabular}{@{}p{0.34\textwidth} p{0.58\textwidth}@{}}
\toprule
\textbf{goal} & \textbf{procedure} \\
\midrule
One transformer, controlled conditions & \texttt{RunTRn\_Base} --- the split is already correct for its size \\
All nine in one pass & \texttt{RunFOR\_Rolling\_PerTrafo\_Async} --- fixed wrapper, size-adaptive split, writes \texttt{\_asyncB} \\
Re-probe after a hardware change & \texttt{RunTR6\_AsyncTune} --- $\sim$1.5\,h, four splits, two periods each \\
\bottomrule
\end{tabular}
\end{center}

\begin{warnbox}
The older \texttt{RunTRn\_Async} wrappers are all hard-coded 6$\times$4, so
\textbf{TR5--TR8 are mis-tuned there} --- prefer \texttt{RunTRn\_Base}.
\texttt{RunTR6\_Async38} and \texttt{RunTR3\_Base4x6} are kept as records of
their experiments; neither split won.
\end{warnbox}

Rules of thumb on this 24-core, \SI{32}{\giga\byte} machine:

\begin{itemize}
  \item keep $K \times \text{threads} = 24$;
  \item 4$\times$6 above $\sim$120 batteries; below that the split is neutral and
        6$\times$4 is the default;
  \item not at full-feeder scale \emph{on this machine}
        (\texttt{ReduceNetwork}~$=0$) --- even the sequential run exhausted its
        \SI{32}{\giga\byte} here. The same run completed elsewhere in
        \SI{82.17}{\hour}, 576/576 optimal, so the wall is local rather than a
        property of the model;
  \item when re-tuning, pick the lowest \emph{wall clock};
  \item one transformer per fresh AIMMS session if the timings are to be compared.
\end{itemize}

Memory observed, peak AIMMS against minimum free: TR4 3.4 / 21.5, TR9 7.2 / 17.3,
TR5 11.2 / 11.8, TR7 14.8 / 8.1, TR6 14.4 / \SI{9.1}{\giga\byte}.

% ================================================================= 8
\section{Files}

\begin{center}
\begin{tabular}{@{}p{0.34\textwidth} p{0.58\textwidth}@{}}
\toprule
\textbf{file} & \textbf{what} \\
\midrule
\texttt{README.md} & this report, in markdown \\
\texttt{timings.csv} & every transformer: size, all three passes, speed-up, inflation, concurrency, validation \\
\texttt{split\_tests.csv} & all three split experiments in one table \\
\texttt{BASELINE\_PROTOCOL.md} & how the controlled base was set up and run \\
\texttt{baseline/} & the campaign's environment log and per-run timings \\
\texttt{probe/} & raw TR6 probe output \\
\texttt{src/} & \LaTeX{} source of this PDF \\
\bottomrule
\end{tabular}
\end{center}

Full FOR results stay at the repository root as
\texttt{FOR\_rolling\_<tag>\_\{gmp,base,async,asyncB\}.csv}, with the
\texttt{\_soc\_} and \texttt{\_fair\_} variants. The model procedures are in
\texttt{MainProject/OPF.ams}, and the snapshots taken before each change are the
\texttt{OPF\_pre\_*.ams.txt} backups beside it.

\end{document}
